\chapter*{Preface}
\markboth{Preface}{Preface}
Welcome to \textit{Microsoft Fabric Master Data Engineer}. This book is a rigorous, hands-on academic guide for learning how modern data engineering changes when storage, compute, governance, orchestration, streaming, machine learning, and business intelligence are delivered through the unified Microsoft Fabric software-as-a-service platform.

The goal is not to memorize product screens. The goal is to develop the architecture judgment expected of a master data engineer: when to use OneLake shortcuts instead of copying data, when to model data as Delta tables versus warehouse tables, how to organize medallion layers, how to reason about capacity units, and how to publish governed analytics products that Power BI users can trust.

\section*{Who This Book Is For}
This book is written for data engineers, analytics engineers, BI developers, cloud architects, database professionals, and technical leaders who need to design and operate data platforms on Microsoft Fabric. It assumes basic SQL literacy and general familiarity with cloud storage, notebooks, and batch pipelines. Python and PySpark are introduced through practical listings, and the exercises gradually move from guided operations to architecture design.

\section*{How This Book Is Organized}
The program is organized as \textbf{6 Parts}, matching the six Fabric milestones of the training plan, and \textbf{24 chapters}, matching 24 instructional weeks plus appendices:
\begin{itemize}
    \item \textbf{Part I: Core Data Infrastructure and OneLake} establishes workspaces, capacities, OneLake, lakehouses, Delta-Parquet, and medallion architecture.
    \item \textbf{Part II: Data Warehousing and Analytics} builds SQL warehouses, semantic models, Direct Lake datasets, and performance-aware reporting paths.
    \item \textbf{Part III: Data Engineering and ETL Orchestration} develops Spark notebooks, Dataflows Gen2, Data Factory pipelines, environments, deployment patterns, and production ETL.
    \item \textbf{Part IV: Real-Time Intelligence and Streaming} introduces Eventstreams, KQL databases, real-time dashboards, and operational telemetry scenarios.
    \item \textbf{Part V: Data Science and Machine Learning} applies Fabric notebooks, feature engineering, MLflow, experiments, and model operationalization to data products.
    \item \textbf{Part VI: Governance, Security, and Exam Prep} covers domains, Purview integration, sensitivity, lineage, access control, monitoring, cost management, and certification-style review.
\end{itemize}

\section*{Reading Paths}
\begin{itemize}
    \item \textbf{New to Fabric:} read Part I in order, complete every hands-on lab, and draw each architecture before implementing it.
    \item \textbf{Experienced Azure data engineer:} focus on OneLake, shortcuts, Direct Lake, capacity management, and Fabric governance differences from separate Azure services.
    \item \textbf{BI-first professional:} read Chapters 1, 5--8, and 21--24 to connect semantic models and Power BI delivery to governed lakehouse design.
    \item \textbf{Platform architect:} read each Friday use case first, then work backward through the week's concepts to validate design trade-offs.
    \item \textbf{Exam preparation:} use each chapter's Key Takeaways, Interview Q\&A, and Exercises as spaced-repetition checkpoints.
\end{itemize}

\section*{Conventions}
Code identifiers, workspace names, lakehouse paths, and commands appear in \texttt{monospace}. PySpark listings use the shared \texttt{pystyle} listing style. Callout boxes flag \textbf{objectives}, \textbf{key terms}, \textbf{definitions}, \textbf{notes}, \textbf{tips}, and \textbf{pitfalls}. Citations point to product documentation and data-management research in the bibliography. Hands-on steps are written for the Microsoft Fabric web experience and PowerShell-friendly Windows workflows unless a cloud service explicitly requires a browser action.
